%% arara directives
% arara: xelatex
% arara: bibtex
% arara: xelatex
% arara: xelatex

%\documentclass{article} % One-column default
\documentclass[twocolumn, switch]{article} % Method A for two-column formatting

\usepackage{preprint}
\usepackage{glossaries}

%% Math packages
\usepackage{amsmath, amsthm, amssymb, amsfonts}

%% Bibliography options
\usepackage[numbers,square]{natbib}
\bibliographystyle{unsrtnat}
%\usepackage{natbib}
%\bibliographystyle{Geology}

%% General packages
\usepackage[utf8]{inputenc}	% allow utf-8 input
\usepackage[T1]{fontenc}	% use 8-bit T1 fonts
\usepackage{xcolor}		% colors for hyperlinks
\usepackage[colorlinks = true,
            linkcolor = purple,
            urlcolor  = blue,
            citecolor = cyan,
            anchorcolor = black]{hyperref}	% Color links to references, figures, etc.
\usepackage{booktabs} 		% professional-quality tables
\usepackage{nicefrac}		% compact symbols for 1/2, etc.
\usepackage{microtype}		% microtypography
\usepackage{lineno}		% Line numbers
\usepackage{float}			% Allows for figures within multicol
%\usepackage{multicol}		% Multiple columns (Method B)

\usepackage{lipsum}		%  Filler text

 %% Special figure caption options
\usepackage{newfloat}
\DeclareFloatingEnvironment[name={Supplementary Figure}]{suppfigure}
\usepackage{sidecap}
\sidecaptionvpos{figure}{c}

% Section title spacing  options
\usepackage{titlesec}
\titlespacing\section{0pt}{12pt plus 3pt minus 3pt}{1pt plus 1pt minus 1pt}
\titlespacing\subsection{0pt}{10pt plus 3pt minus 3pt}{1pt plus 1pt minus 1pt}
\titlespacing\subsection{0pt}{8pt plus 3pt minus 3pt}{1pt plus 1pt minus 1pt}

% ORCiD insertion
\usepackage{tikz,xcolor,hyperref}

\definecolor{lime}{HTML}{A6CE39}
\DeclareRobustCommand{\orcidicon}{
	\begin{tikzpicture}
	\draw[lime, fill=lime] (0,0) 
	circle [radius=0.16] 
	node[white] {{\fontfamily{qag}\selectfont \tiny ID}};
	\draw[white, fill=white] (-0.0625,0.095) 
	circle [radius=0.007];
	\end{tikzpicture}
	\hspace{-2mm}
}
\foreach \x in {A, ..., Z}{\expandafter\xdef\csname orcid\x\endcsname{\noexpand\href{https://orcid.org/\csname orcidauthor\x\endcsname}
			{\noexpand\orcidicon}}
}
% Define the ORCID iD command for each author separately. Here done for two authors.
\newcommand{\orcidauthorA}{0009-0004-5890-4819}
\newcommand{\orcidauthorB}{0009-0006-4869-9705}

%%%%%%%%%%%%%%%%   Title   %%%%%%%%%%%%%%%%
\title{Heuristic Solutions for the Trigger Arc Traveling Salesman Problem: An Application of GRASP}

% Add watermark with submission status
\usepackage{xwatermark}
% Left watermark
\newwatermark[firstpage,color=gray!60,angle=90,scale=0.32, xpos=-4.05in,ypos=0]{\href{https://doi.org/}{\color{gray}{Publication doi}}}
% Right watermark
\newwatermark[firstpage,color=gray!60,angle=90,scale=0.32, xpos=3.9in,ypos=0]{\href{https://doi.org/}{\color{gray}{Preprint doi}}}
% Bottom watermark
\newwatermark[firstpage,color=gray!90,angle=0,scale=0.28, xpos=0in,ypos=-5in]{*correspondence: \texttt{joan.salvasoler@h2o.ai}}

%%%%%%%%%%%%%%%  Author list  %%%%%%%%%%%%%%%
\usepackage{authblk}
\renewcommand*{\Authfont}{\bfseries}
\author[1]{Joan Salva Soler\orcidA{}}
\author[2]{Grégoire de Lambertye\orcidB{}}

\affil[1]{H2O.ai, TU Wien}
\affil[2]{TU Wien}
 
%%%%%%%%%%%%%%    Front matter    %%%%%%%%%%%%%%
\begin{document}

\twocolumn[ % Method A for two-column formatting
  \begin{@twocolumnfalse} % Method A for two-column formatting
  
\maketitle

\begin{abstract}
 purus elit, vestibulum ut, placerat ac, adipiscing vitae, felis.
 Curabitur dictum gravida mauris. Nam arcu libero, nonummy
 eget, consectetuer id, vulputate a, magna. Donec vehicula au
gue eu neque. Pellentesque habitant morbi tristique senectus
 et netus et malesuada fames ac turpis egestas. Mauris ut leo.
 Cras viverra metus rhoncus sem. Nulla et lectus vestibulum
 urna fringilla ultrices. Phasellus eu tellus sit amet tortor gravida
 placerat. Integer sapien est, iaculis in, pretium quis, viverra ac,
 nunc. Praesent eget sem vel leo ultrices bibendum. Aenean

\end{abstract}

\keywords{TSP \and TA TSP \and GRASP \and Local Search \and Trigger Arc TSP \and Optimisation} % (optional)

\vspace{0.35cm}

  \end{@twocolumnfalse} % Method A for two-column formatting
] % Method A for two-column formatting

%\begin{multicols}{2} % Method B for two-column formatting (doesn't play well with line numbers), comment out if using method A


%%%%%%%%%%%%%%%  Main text   %%%%%%%%%%%%%%%
% \linenumbers
\newacronym{tsp}{TSP}{Traveling Salesman Problem}
\newacronym{tatsp}{TA-TSP}{Trigger Arc TSP}
\newacronym{grasp}{GRASP}{Greedy Randomized Adaptive Search Procedures}

\section{Introduction}
\lipsum[1]
\gls{tsp}

\section{Background}
\label{sec:background}

\subsection{The TA TSP}

The \gls{tatsp}, is a variation of the \gls{tsp} introduced by~\cite{Cerrone}. 
Unlike traditional TSP formulations, which typically assume static or purely time-dependent arc costs, 
the \gls{tatsp} incorporates a dynamic cost mechanism that directly reflects the sequence of traversed arcs within a tour.

Specifically, this problem introduces a set of Relations, comprising two distinct arc types:
\begin{itemize}
\item Trigger Arcs: These are designated arcs whose traversal can influence the cost of other arcs within the network.
\item Target Arcs: These are arcs whose traversal cost is subject to dynamic modification.
\end{itemize}
Each relation defines a specific cost that a target arc will assume if its corresponding trigger arc is the last trigger 
arc visited immediately preceding the target arc's traversal. This mechanism is critical, as it establishes a pronounced 
path-dependency, making the total tour cost contingent on the specific sequence of arcs traversed.

A formal definition of the problem is given by~\cite{Cerrone} as:
consider a directed graph $G = (N,A)$ with weights on the arcs. 
The node $0 \in N$ is designated as the starting node (depot). Let $c(a) \in \mathbb{R}^+ \forall a \in A$ 
be defined as the cost of traversing the arc $a$. For each arc $a=(h, k)$, a set of relations $R_a = \{(a_1, a)| a_1 \in A\}$ 
is associated. Finally, let $c(r) \in \mathbb{R}^+ \forall r \in R_a$, be the traversal cost of the arc $a$ if the relation $r$ 
is active. Let $T = (a_1, a_2, \ldots, a_{|N|})$ be the ordered sequence of arcs starting at node $0$ representing a Hamiltonian 
cycle in $G$. The relation $r = (a_i, a_j)$ is active if and only if the arcs $a_i, a_j \in T$, and the arc $a_i$ precedes the arc 
$a_j$ in $T$ and there is no active relation $r_1 = (a_k, a_j) \in R_{a_j}$ such that $a_k$ precedes $a_i$ in $T$. It follows that 
for each arc $a$, at most one relation can be active. As a result, the traversal cost of the arc $a = (h, k)$ will be equal to $c(a)$ 
if there are no active relations in $R_a$ or $c(r)$ if $r$ is the only active relation in $R_a$.

This particular problem formulation is designed to facilitate the modeling of complex real-world scenarios where the act of 
traversing an arc by a vehicle or agent can intrinsically alter the properties or costs associated with subsequent segments of 
the underlying road network or system state. Building upon its initial definition, this work further investigates the \gls{tatsp}, 
aiming to explore its properties and develop methods for its resolution.

\subsection{The GRASP approach}

The solution developped in this paper mainly rely on a proven metaheuristic introduced by Feo et al.\cite{Feo1995}, the \gls{grasp}. 
This powerful and flexible framework is well-suited for tackling complex combinatorial optimization problems, 
including those exhibiting the intricate path-dependencies seen in the \gls{tatsp}. It operates through an iterative 
two-phase process: a construction phase and a local search phase. 

The construction phase leverages a greedy randomized strategy to build a feasible solution. This involves creating diverse feasible 
solutions at minor computationnal costs. The diversity of the solution is carried on by a degree of randomness. 

Following the construction of a solution, the local search phase then systematically explores its neighborhood to improve it, 
typically by applying moves that lead to better objective function values. The best solution found across multiple iterations of 
these two phases is ultimately returned. 

This rigorous yet adaptive approach makes \gls{grasp} particularly well-suited for problems where the optimal solution is difficult 
to find through exact methods.

\section{Methodology}
Overview of the approach 

\subsection{Constructing initial solution}
\lipsum[5]

\subsection{Local search}
\lipsum[6]

\subsubsection{Delta evaluation}
\lipsum[7]

\section{Experimental Setup and Data}
\lipsum[8]

\subsection{Datasets}
The method described earlier has been used over 3 datasets 

\subsection{Experimental Setup}
\lipsum[10]

\section{Results and Discussion}
\lipsum[11]

\section{Conclusion}
\lipsum[12] \cite{Cerrone}

%%%%%%%%%%%% Supplementary Methods %%%%%%%%%%%%
%\footnotesize
%\section*{Methods}

%%%%%%%%%%%%% Acknowledgements %%%%%%%%%%%%%
%\footnotesize
%\section*{Acknowledgements}

%%%%%%%%%%%%%%   Bibliography   %%%%%%%%%%%%%%
\normalsize
\bibliography{references}

%%%%%%%%%%%%  Supplementary Figures  %%%%%%%%%%%%
%\clearpage

%%%%%%%%%%%%%%%%   End   %%%%%%%%%%%%%%%%
%\end{multicols}  % Method B for two-column formatting (doesn't play well with line numbers), comment out if using method A
\end{document}